\documentclass[10pt]{article}

\usepackage[utf8]{inputenc}

\usepackage[T1]{fontenc}

\usepackage{amsmath}

\usepackage{amsfonts}

\usepackage{amssymb}

\usepackage[version=4]{mhchem}

\usepackage{stmaryrd}

\usepackage{graphicx}

\usepackage[export]{adjustbox}

\graphicspath{ {./images/} }



\begin{document}

но. Говорят, что схема (3) аппроксимирует задачу (1), (2) с $m$-м порядком аппроксимации, өсли



\[

\|\psi\|_{(2 h)}=O\left(|h|^{m}\right),\|v\|_{\left(3_{h}\right)}=O\left(|h|^{m}\right), m>0 .

\]



Схема (3) имеет $m$-й порядок точности или сходится со скоростью $O(|h| m)$, если



\[

\left\|y_{h}-u_{h}\right\|_{\left(1_{h}\right)}=O\left(|h|^{m}\right) \text {. }

\]



Для сходимости разностной схемы одной только аппроксимации, вообще говоря, недостаточно: надо потребовать еще, чтобы разностная задача (3) была корректна. Именно, справедливо следующее утверждение: если разностная схема (3) корректна и имеет $m$-й порядок аппроксимации, то она сходится со скоростью $O\left(|h|^{m}\right)$ (cm. [28]).



Возможны и другие подходы к построению Р. с. т. Так, в теории Лакса (см. [8]) сходимость разностных схем изучается не в пространствах сеточных функций, а в пространстве решений исходной дифферевциальной задачи; здесь доказана т. н. теорема әквивалентности: если исходная задача (1), (2) корректна и схема (3) аппроксимирует задачу (1), (2), то устойчивость необходима и достаточна для сходимости. Возможны и другие постановки вопроса о связи устойчивости и сходимости (см., напр., [9]). Исследование сходимости разностных схем проводилось и в пространствах обобщенных репений (см. [10]).



Требования к разностным схемам. При расчетах на современных ЭВМ не всегда достаточно требовать от разностной схемы только сходимости при $|h| \rightarrow 0$. Использование реальных сеток при конечном шаге сетки предъявляет к схемам ряд дополнительных требований. Они сводятся к тому, что разностная схема помимо обычной ашпроксимации и устойчивости должна хорошо моделировать характерные свойства исходного диффференциального уравнения. Кроме того, разностная схема должна удовлетворять определенным условиям простоты реализации вычислительного алгоритма. Ниже рассматриваются нек-рые из әтих дополнительных требований.



Под одно родно й $\mathrm{p}$ а зностно й с хемой (см. [1], [11]) понимается разностная схема, вид к-рой не зависит ни от выбора конкретной задачи из данного класса, ни от выбора разностной сетки. Во всех узлах сетки для любой задачи из данного класса разностные уравнения имеют один и тот же вид. К однородным разностным схемам относятся, в частности, схемы сквозного счета для решения уравнений с сильно меняющимися или разрывными коэффициентами. Схемы сквовного счета не предусматривают явного выделения точек разрыва коэфффициентов и позволяют вести выqисления по одним и тем же формулам. Широко используются схөмы сквозного счета при расчетах разрывных решений уравнений газовой динамики (см. [1], [5], [12]).



Требование к о н с е в а т и в н о с т и р а н о с т н й с х е м ы означает, что данная разностная схема имеет на сетке такой же закон сохранения, что и исходное дифференциальное уравнение. В частности, если $L$ - самосопряженный оператор и схема (3) консервативна, то $L_{h}$ - самосопряженный в $H_{h}$ оператор, т. е. консервативная схөма сохраняет свойство самосопряженности.



Регулярным методом получения консервативных однородных схем является т. н. м е т о д б а л а н с а, или ин те г ро - и т е р по ляци н ны й мето д. Сущность метода баланса состоит в аппроксимации на разностной сетке интегрального закона сохранения (уравнения баланса), соответствуюшцео данному дифференгиајьному уравнениіо. Метод б́аланса нашелі широкое применение при аппроксимации уравнений с переменными, в том числе п разрывными, коәфффициентами. Другой группой методов построения разностных схем, сохраняющих свойства самосопряженности и положительности исходного оператора, являются методы, основанные на вариационных принципах (метод Ритца, метод конечных элементов) (см. Разностная вариационнал схема).



Прп конструировании разностных схем Для уравнений газовой динамики нашел примөнөние принцип полной консервативности (см. [5]). Для уравнөний гиперболич. типа оказалось полезным исследование дисперсиянных свойств соответствующих разностных уравнений (см. [13]).



Если известно, что при $t \rightarrow \infty$ решениө исходной дифффренциальной задачи стремится к нулю, то естественно требовать того же и от решения аппроксимирующей разностной задачи. Схемы, обладающие әтим свойством, наз. а с и м І т о т Ич е с к И у с то йч И в Ы М И (см. [4]).



Другой подход к построению разностных схем лучшего качества состоит в получөнии схөм, удовлетворяющих тем же априорным оценкам, к-рые характерны (неулучшаемы) для исходных диффферөнциальных уравнений (cм. [14]).



Экономичные разностные схемы для мдогомерных задач математической физики. При решенид систем разностных уравнөний, аппроксимирующих нестационарные многомерные задачи матөматич. физики (с числом пространственных пөременных два или болеө), возникают специфич. затруднения, связанные с тем, что число арифметич. операций, нөобходимых для отыскания решения на новом врөмөнном слоө, резко возрастает при измельчөнии сөтки. В то жө время решение одномерных вестационарных задач осущөствЈяется методом прогонки, к-рый әкономпчен в том смысле, что он требует конечного (не зависящего от шага сетки $h$ ) числа действий на одну точку сетки.



\includegraphics[max width=\textwidth, center]{2023_07_08_9af9c8fa6e2b1f3c0653g-1}

н о й, если отношение числа дөйствий, необходимых для отыскания решения на новом временно́м слое к числу узлов пространственной сетки, нө зависпт от числа узлов сетки. Обычныө неявные равностные схемы не являются экономичными. Напболөө әффективным приемом построения әкономичных равностных схем является метод сведения многомөрных задач к нескольким одномерным задачам (мөтод переменвых направлений, метод расщепления) (см. [1], [15], [16]).



Новые алгоритмы стимулировали и новый подход к основным понятиям теории разностных схем: аппроксимации, устойчивости, сходимости. Напр., оказалось плодотворным понятие суммарвой аппрокспмации или аппроксимации в слабом смысле (см. [1], [16]). Сформулирован принцип аддитивности, позволяющий в общем случае строцть экономичные разностные схемы, обладающие свойством суммарной аппроксимации (см. [17]). Дальнейшим обобщением разностных схем переменных направлений явплись схемы с уравнениями на графах и векторные схемы (см. [18]).



Методы исследования корректности и сходимости разностных схем. Для линейных задач из устойчивости и аппроксимации следует сходимость. Поәтому основное внимание в Р. с. т. уделяется получению априорных оценок, из к-рых следует корректность задач в той или иной норме. Методы получения априорных оценок для разностных схем во многом аналогичны тем же метөдам в теории дифффөренциальных уравнений, напр. можно указать следующие мөтоды: разделение переменных, преобразование Фурье, принцип максимума, энергетич. неравенства.



Методы решения сеточных уравнений. Јюбой сеточјый метод для диффференциальных уравнений приводит к большим системам линейных алгебраич. уравнений.





\end{document}